\emph{Preliminaries.} Write $\tau = \langle\cdot\rangle_{0}$. By Theorem \claimref{synthesis.signed-directions-and-the-integer-clifford-algebra} the blades are a basis and each blade $b$ has $b^{-1} = \pm b$, so $b^{\dagger} = b$ exactly when $b^{2} = 1$ and $b^{\dagger} = -b$ otherwise; $\dagger$ is an anti-automorphism, and $\tau(x^{\dagger}x)$ is the sum of the squares of the coefficients of $x$ (clause (iii) of Theorem \claimref{synthesis.chsh-bound-for-two-subsystems-of-a-clifford-algebra-by-a-rational-certificate}). The cone $\mathcal{K}$ is closed under sums and under multiplication by non-negative elements of $k$, and contains $1 = 1 \cdot 1^{\dagger}$ and every projector $E = EE^{\dagger}$; if $E$ is a projector so is $1 - E$, so projectors are effects. Two remarks used below. (A) If $E$ is an effect and $0 \le t \le 1$ in $k$, then $tE$ is an effect, since $1 - tE = t(1 - E) + (1 - t)$. (B) If $c, c' \in \mathcal{K}$ and $c + c'$ is an effect, then $c$ and $c'$ are effects, since $1 - c = (1 - c - c') + c'$.

\emph{(i).} The self-adjoint elements are the span of the self-adjoint blades because $\dagger$ multiplies each blade by $\pm 1$. For a blade $b$ of square $1$, $(1 \pm b)^{\dagger} = 1 \pm b$ and $(1 \pm b)^{2} = 2(1 \pm b)$. For $s = \sum_{b}c_{b}b$ (the scalar blade included), $n_{s} - s = \sum_{b}|c_{b}|\,(1 - \operatorname{sgn}(c_{b})\,b)$, and each term is $|c_{b}|$ times an element $\tfrac12(1 \mp b)(1 \mp b)^{\dagger}$ (for the scalar blade the term is $0$ or $2|c_{b}|$); likewise $n_{s} + s$. So $s = \tfrac12(n_{s} + s) - \tfrac12(n_{s} - s)$ is a difference of elements of $\mathcal{K}$, and for $s \ne 0$ the element $(n_{s} + s)/(2n_{s})$ is an effect, with complement $(n_{s} - s)/(2n_{s})$, so $s$ is a combination of that effect and of the effect $1$.

\emph{(ii).} A state restricts to a weight: non-negativity on $\mathcal{K}$ is the positivity clause with linearity, and additivity is linearity. Conversely let $\mu$ be a weight. Step 1: for an effect $E$ and integers $0 \le j \le n$, $n \ge 1$, the elements $jE/n$ are effects by (A), so additivity gives $\mu(jE/n) = j\mu(E/n)$ and $\mu(E) = n\mu(E/n)$: $\mu(rE) = r\mu(E)$ for rational $r \in [0,1]$. Step 2, used only when $k \ne \Q$: for $t \in k$ with $0 \le t \le 1$ and rationals $r \le t \le r'$ in $[0,1]$, $\mu(tE) = \mu(rE) + \mu((t - r)E) \ge r\mu(E)$ and $\mu(r'E) = \mu(tE) + \mu((r' - t)E) \ge \mu(tE)$, by additivity and non-negativity; so $\mu(tE)$ and $t\mu(E)$ both lie between $r\mu(E)$ and $r'\mu(E)$, and since $k$ is Archimedean $r' - r$ can be made smaller than any positive element, whence $\mu(tE) = t\mu(E)$. Step 3: for $c \in \mathcal{K}$ there is an integer $n \ge 1$ with $c/n$ an effect: by (i) $n_{c} - c \in \mathcal{K}$, and any integer $n \ge n_{c}$ (which exists as $k$ is Archimedean) has $1 - c/n = \bigl((n_{c} - c) + (n - n_{c})\bigr)/n \in \mathcal{K}$. Put $\omega(c) = n\mu(c/n)$; by Step 1 this does not depend on $n$ (compare $n$ and $n'$ through $nn'$). For $c, c' \in \mathcal{K}$ choose $n$ with $(c + c')/n$ an effect; by (B) $c/n$ and $c'/n$ are effects and additivity of $\mu$ gives $\omega(c + c') = \omega(c) + \omega(c')$. By Steps 1 and 2, $\omega(tc) = t\omega(c)$ for $t \ge 0$ in $k$, and $\omega \ge 0$ on $\mathcal{K}$. Step 4: by (i) every self-adjoint $s$ is $c - c'$ with $c, c' \in \mathcal{K}$; put $\omega(s) = \omega(c) - \omega(c')$, well defined because $c - c' = d - d'$ gives $c + d' = d + c'$ and additivity on $\mathcal{K}$ applies. This $\omega$ is $k$-linear on the self-adjoint elements. Step 5: put $\omega(x) = \omega(\tfrac12(x + x^{\dagger}))$ for all $x$. Then $\omega$ is $k$-linear, $\omega(x^{\dagger}) = \omega(x)$, $\omega(xx^{\dagger}) \ge 0$ because $xx^{\dagger} \in \mathcal{K}$, $\omega(1) = \mu(1) = 1$, and $\omega = \mu$ on effects (take $n = 1$): $\omega$ is a state in the sense of Premise \claimref{synthesis.premise-positivity-of-values}. Uniqueness: two states that agree on effects agree on their span, the self-adjoint elements, and the reality clause gives $\omega(x) = \omega(\tfrac12(x + x^{\dagger}))$. For $k = \Q$ Step 2 is not needed, and Steps 1, 3, 4, 5 use only additivity: they give a $\Q$-linear real $\omega$ for any additive map on effects with values in $\Q$; then $\omega(xx^{\dagger}) \ge 0$ for all $x$ exactly when $\mu \ge 0$ on effects, because every element of $\mathcal{K}$ is a positive integer multiple of an effect.

\emph{(iii).} $\tfrac{3}{10}\Pi_{e} + \tfrac{4}{10}\Pi_{p} = \tfrac{7}{20} + \tfrac{3}{20}e_{1} + \tfrac{4}{20}p_{1} = \tfrac12E_{a} + \tfrac{1}{10}$. Each term is a multiple of a projector or of $1$ by a number in $[0,1]$, hence an effect by (A), and the common value is an effect because its complement is $\tfrac12(1 - E_{a}) + \tfrac{4}{10} \in \mathcal{K}$; so by (B) the partial sums are effects and additivity applies to both sides: $\tfrac12\mu(E_{a}) + \tfrac{1}{10} = \tfrac{3}{10}\mu(\Pi_{e}) + \tfrac{4}{10}\mu(\Pi_{p})$, using Step 1 of (ii) for the rational multiples. With the two given values, $\mu(E_{a}) = \tfrac{9}{10}$. Here $\Cl_{k}(1,1)$ is one mode in the sense of Proposition \claimref{tower.ladder-pairs-satisfy-car-vectors}, with generators $e_{1}$, $f_{1}$ of squares $1$, $-1$. Its self-adjoint elements are the span of $1$, $e_{1}$, $p_{1}$ (by (i): $f_{1}^{2} = -1$), and $\tfrac12(1 + n)$ with $n = xe_{1} + zp_{1}$ is idempotent exactly when $n^{2} = x^{2} + z^{2} = 1$ (the elements $e_{1}$, $p_{1}$ anticommute), while an idempotent $c + \tfrac12 n'$ with $c \ne \tfrac12$ forces $n' = 0$ and $c \in \{0, 1\}$; two such projectors $\tfrac12(1 + n)$, $\tfrac12(1 + n')$ have product zero only if $n' = -n$. So the families of orthogonal nonzero projectors are the pairs $\{E, 1 - E\}$, and a function $g$ on these projectors is additive on them exactly when $g(\tfrac12(1 - n)) = 1 - g(\tfrac12(1 + n))$, which holds for $\tfrac12(1 + z^{3})$ since $z \mapsto z^{3}$ is odd; it is non-negative because $|z| \le 1$. Its values $1$, $\tfrac12$, $\tfrac{189}{250}$ are the script's, and $\tfrac12 \cdot \tfrac{189}{250} + \tfrac{1}{10} = \tfrac{239}{500} \ne \tfrac{275}{500}$.

\emph{(iv).} $\Gamma^{\dagger} = \Gamma$ gives $(\Gamma x\Gamma)^{\dagger} = \Gamma x^{\dagger}\Gamma$, so $E_{\Gamma}$ commutes with $\dagger$; and $\Gamma yy^{\dagger}\Gamma = (\Gamma y)(\Gamma y)^{\dagger}$. Hence $E_{\Gamma}(\mathcal{K}) \subset \mathcal{K}$, every element of $E_{\Gamma}(\mathcal{K})$ is fixed because $E_{\Gamma}$ is idempotent (Theorem \claimref{singularity.descent-parity-average} for one mode; in general $\Gamma^{2} = 1$ gives $E_{\Gamma}E_{\Gamma} = E_{\Gamma}$), and a fixed element $c$ of $\mathcal{K}$ is $E_{\Gamma}(c)$. The statement about $\varphi \circ E_{\Gamma}$ follows. In the example $1 + 2p_{1}$ is self-adjoint and commutes with the total parity $p_{1}p_{2}$ of Theorem \claimref{synthesis.mode-split-of-the-clifford-algebra-and-its-klein-twist-to-the-tensor-product}, so $\varphi$ is real and invariant; $\varphi(\tfrac12(1 - p_{1})) = \tfrac12(1 - 2) = -\tfrac12$, and $\tfrac12(1 - p_{1})$ is a projector. Rider: the averages over the commuting self-adjoint $p_{j}$ commute, and their composition $E$ annihilates every blade that anticommutes with some $p_{j}$ and fixes the others, which are, up to sign, the products of the $p_{j}$; their span is the span of the $2^{\ell}$ orthogonal projectors $\Pi_{\chi} = \prod_{j}\tfrac12(1 + \chi_{j}p_{j})$. A fixed element of $\mathcal{K}$ is $c = \sum_{\chi}c_{\chi}\Pi_{\chi}$ with $c_{\chi}\tau(\Pi_{\chi}) = \tau(\Pi_{\chi}c) \ge 0$, because $\tau(\Pi_{\chi}yy^{\dagger}) = \tau((\Pi_{\chi}y)(\Pi_{\chi}y)^{\dagger}) \ge 0$; so an invariant functional is non-negative on $\mathcal{K}$ exactly when it is non-negative on the $\Pi_{\chi}$. For the grading: conjugation by the total parity (or, in general, the grading automorphism) fixes even and negates odd elements, so for $y = y_{0} + y_{1}$ split into even and odd parts $E_{\Gamma}(yy^{\dagger}) = y_{0}y_{0}^{\dagger} + y_{1}y_{1}^{\dagger}$, the cross terms being odd. For a generator $g$, $gg^{\dagger} = gg^{-1} = 1$, so $y_{1}y_{1}^{\dagger} = (y_{1}g)(y_{1}g)^{\dagger}$ with $y_{1}g$ even. So the fixed elements of $\mathcal{K}$ are combinations, with non-negative coefficients, of hermitian squares of even elements, and conversely such combinations are fixed. A functional vanishing on odd elements is invariant, so it is non-negative on $\mathcal{K}$ exactly when it is non-negative on hermitian squares of even elements; normalisation and reality restrict in the obvious way.

\emph{(v).} With $p_{1} = f_{1}e_{1}$: $e_{1}\,e_{2}p_{1} = e_{1}e_{2}f_{1}e_{1} = e_{2}f_{1} = u$ (moving the last $e_{1}$ to the front costs two transpositions), $p_{1}p_{2} = \Gamma$, and the two mixed products $e_{1}p_{2}$ and $p_{1}\,e_{2}p_{1} = e_{2}$ are odd. So $a_{i}b_{j} = x_{i}x'_{j}u + z_{i}z'_{j}\Gamma + (\text{odd})$ and $E_{\Gamma}(a_{i}b_{j}) = x_{i}x'_{j}u + z_{i}z'_{j}\Gamma$; summing with the CHSH signs gives $E_{\Gamma}(C) = Xu + Z\Gamma$. $u^{2} = e_{2}f_{1}e_{2}f_{1} = -e_{2}^{2}f_{1}^{2} = 1$, so $u$ is a blade of square $1$, self-adjoint, and it commutes with $\Gamma$ because it is even. $E_{\Gamma}(a_{i}) = z_{i}p_{1}$ and $E_{\Gamma}(b_{j}) = z'_{j}p_{2}$ ($e_{1}$ and $e_{2}p_{1}$ are odd), whose CHSH element is $Z\Gamma$. For an even state, $\omega(C) = \omega(E_{\Gamma}(C)) = X\omega(u) + Z\omega(\Gamma)$. For self-adjoint $w$ with $w^{2} = 1$, $1 \pm w = \tfrac12(1 \pm w)(1 \pm w)^{\dagger}$ and positivity give $|\omega(w)| \le 1$; hence $|\omega(C)| \le |X| + |Z|$. The four elements $\tfrac14(1 + \epsilon u)(1 + \epsilon'\Gamma)$, $\epsilon, \epsilon' = \pm 1$, are projectors (products of commuting projectors), even, of scalar part $\tfrac14$ because $u$, $\Gamma$ and $u\Gamma$ are non-scalar blades; the state $x \mapsto 4\langle\tfrac14(1 + \epsilon u)(1 + \epsilon'\Gamma)x\rangle_{0}$ is even and has $\omega(u) = \epsilon$, $\omega(\Gamma) = \epsilon'$, so $\omega(C) = \epsilon X + \epsilon'Z$, which is $|X| + |Z|$ for suitable signs. The bound: with the plane vectors $\alpha_{i} = (x_{i}, z_{i})$, $\beta_{j} = (x'_{j}, z'_{j})$ of unit length, $X + Z = \alpha_{1} \cdot (\beta_{1} + \beta_{2}) + \alpha_{2} \cdot (\beta_{1} - \beta_{2})$, so $(X + Z)^{2} \le 2\bigl(|\beta_{1} + \beta_{2}|^{2} + |\beta_{1} - \beta_{2}|^{2}\bigr) = 8$ by the Cauchy-Schwarz inequality in $k^{2}$; replacing $z'_{j}$ by $-z'_{j}$ gives the same for $X - Z$, and $|X| + |Z|$ is $|X + Z|$ or $|X - Z|$. The numbers at the stated directions, $X = 1 \cdot (\tfrac35 + \tfrac45) = \tfrac75$ and $Z = 1 \cdot (\tfrac45 + \tfrac35) = \tfrac75$, and the value $\tfrac{14}{5}$, are confirmed by the script.
